It also follows that the first statement in 6.2 holds without supposing $G$ locally of finite type over $S$, provided the immersion $K\to G$ is of finite presentation.

\numpara{6.5}
As announced in 5.6, theorem 6.2 allows one to extend to group preschemes satisfying a') and b) above certain results established by another method and under more restrictive conditions for groups of multiplicative type. This is so for results 5.2, the beginning of 5.3, 5.4 and the bis variants of the preceding results 5.9 and 5.10. It is likewise so for the results of IX, Nos.~5 and 6, excluding IX~6.8 (already false for a homomorphism of abelian schemes $u:H\to G$ over the spectrum $S$ of a discrete valuation ring with residue characteristic $p>0$: it can indeed happen that $\Ker u$ has as generic fiber the identity group, and as special fiber a radicial group not reduced to the identity group).

\numpara{6.6}
We have just given an example of a rigidity property for groups of multiplicative type that is not shared by abelian schemes. Another example, extremely important, is the fact that the existence theorem for infinitesimal extensions of homomorphisms IX~3.6 no longer holds when $H$ is an abelian scheme. Thus, an abelian scheme $\ne0$ over a field admits nontrivial infinitesimal variations, unlike what happens for a group of multiplicative type---which is the infinitesimal aspect of the fact that there exists a ``moduli theory'' (moreover far from complete) for abelian varieties, whereas moduli theory for groups of multiplicative type is empty. Another ``global'' aspect of this infinitesimal difference is that if $H$ is an abelian scheme over a locally noetherian $S$, and $G$ a commutative $S$-group prescheme locally of finite type over $S$, then one can show that $\uHom_{S\text{-gr}}(H,G)$ is representable by a prescheme locally of finite type over $S$, but, unlike what happens for $H$ of multiplicative type, this prescheme \emph{is not étale} over $S$, but only unramified over $S$. Thus, if $S$ is, for example, the spectrum of a complete discrete valuation ring, $H$ and $G$ abelian schemes over $S$, there can exist homomorphisms $H_0\to G_0$ on the special fibers that do not come ``by specialization'' from a homomorphism on the generic fibers.

\numpara{6.7}
Theorem 6.2 applies whenever $H$ is an abelian scheme over $S$, or more generally an extension of such a scheme by a torus. Indeed, since the question is local on $S$, one can suppose that there exists a prime number $\ell$ prime to the residue characteristics of $S$, and one sees that it then suffices to take for $E$ the set of powers of $\ell$ to satisfy conditions a') and b).

An argument analogous to that of 6.1 gives us the
\begin{theorem}{6.8}\label{XI.6.8}
Let $X$ be a prescheme smooth over $S$, with nonempty connected geometric fibers. Then for every closed sub-prescheme $Y$ of $X$, the functor $\prod_{X/S}Y/X$ is representable by a closed sub-prescheme $T$ of $S$. If $Y$ is of finite presentation over $X$, then $T$ is of finite presentation over $S$.
\end{theorem}
Since $f:X\to S$ is faithfully flat locally of finite presentation, it is covering for the fpqc topology. Since, on the other hand, $T=\prod_{X/S}Y/X$ is obviously a subsheaf of $S$ (for the fpqc topology), it follows that the question of representability of $T$ by a closed sub-prescheme of $S$ is of a local nature on $S$ for the fpqc topology, and the same is true of the question of deciding whether $T$ is of finite presentation over $S$. On making then the base change $S'\to S$, with $S'=X$, one is reduced to the case where $X$ admits a section $e$ over $S$. One can in addition suppose $S$ affine and a fortiori quasi-compact. One then has:
\begin{corollary}{6.9}\label{XI.6.9}
Under the conditions of 6.8, suppose that $S$ is quasi-compact and that $X$ admits a section $e$ over $S$. For every integer $n\ge0$, let $X_n$ be the sub-prescheme of $X$ that is the infinitesimal neighborhood of order $n$ of the section $e$. Suppose $Y$ of finite presentation over $X$. Then there exists an integer $n\ge0$ such that one has $\prod_{X/S}Y/X=\prod_{X_n/S}Y_n/X_n$ (where $Y_n=Y\cap X_n$).
\end{corollary}
This corollary does indeed imply 6.8 when $Y$ is of finite presentation over $X$, thanks to VIII~6.4, for, since $X$ is smooth over $S$, $X_n$ is finite and locally free over $S$ and a fortiori is ``essentially free'' over $S$ in the sense of VIII~6.1, hence $\prod_{X_n/S}Y_n/X_n$ is representable by a closed sub-prescheme of $S$. In addition, the proof in \loccit, or the reduction to the noetherian case, shows us immediately that the said closed sub-prescheme of $S$ is of finite presentation over $S$.

Let us first prove 6.9, hence 6.8, when $S$ is noetherian. Let $T_n=\prod_{X_n/S}Y_n/X_n$; then the $T_n$ form a decreasing sequence of closed sub-preschemes of $S$, and since $S$ is noetherian, this sequence is stationary. Let $R=\bigcap_{n\ge0}\prod_{X_n/S}Y_n/X_n$ be their common value for large $n$; one obviously has $T\subset R$, and it suffices to establish that one has $R\subset T$. On making the base change $R\to S$, one is reduced to the case where $R=S$, \ie $Y_n=X_n$ for every $n$, \ie $Y\supset X_n$ for every $n$, and to proving then $T=S$, \ie $Y=X$. Now $Y\supset X_n$ for every $n$ implies (thanks to the fact that $X$ is locally noetherian) that $Y$ is, near every point of $e(S)$, an induced open sub-prescheme of $X$, so there exists an induced open subset $U$ of $X$, containing $e(S)$, such that $U\subset Y$. By IX~4.3, since the fibers of $X/S$ are integral, $U$ is schematically dense in $X$, hence ($Y$ being a closed sub-prescheme containing $U$) one has $Y=X$. This proves 6.9, hence 6.8, in this case.

The general case proceeds by reduction to the preceding case. For every $s\in S$, there exists an affine open neighborhood $U$ of $s$ and an affine open neighborhood $V$ of $e(s)$ such that $f(V)\subset U$. Then $f(V)$ is an open neighborhood of $s$ contained in $U$, and if $S'$ is an affine open neighborhood of $s$ contained in $e^{-1}(V)\cap f(V)$, and $X'=V\cap f^{-1}(S')$, then $X'$ and $S'$ are affine open subsets of $X$, resp. $S$, and $X'/S'$ admits a section. Because of the local nature of 6.8 and 6.9, one can suppose $S=S'$. I say that one then has $\prod_{X/S}Y/X=\prod_{X'/S}Y'/X'$, where $Y'=Y\cap X'$; indeed, by IX~4.6, $X'$ is schematically dense in $X$ (at least when $X$ is quasi-separated over $S$, hence $X'$ is retrocompact in $X$; but in fact one can show without difficulty that IX~4.6 remains valid without the retrocompactness hypothesis), and likewise for every base change $S_1\to S$, $X'\times_S S_1$ is schematically dense in $X\times_S S_1$, whence at once the asserted equality. This reduces us to the case where $X=X'$, so one can suppose $S$ and $X$ affine.

In addition, if $X=\Spec(B)$ and if $J$ is the ideal of $B$ defining $Y$, $J$ is the inductive limit of its subideals of finite type, so $Y$ is an intersection of closed sub-preschemes $Y_i$ of $X$ that are of finite presentation over $S$, and consequently $\prod_{X/S}Y/X=\bigcap_i\prod_{X/S}Y_i/X$, which reduces us, to prove 6.8, to the case where $Y$ is of finite presentation over $X$. It then suffices to prove 6.9 with $S$ and $X$ affine. But then $X$ and $Y$ over $S$ come by the base change $S\to S_0$ from an analogous situation $X_0$ and $Y_0$ over $S_0$, with $S_0$ noetherian, which reduces us to the case where $S$ is noetherian, which has already been treated. This completes the proof of 6.8 and 6.9.
% typo: "S et noethérien" -> "S est noethérien".
% typo?: Y_i of finite presentation over S, rather than X, kept as printed.

\begin{corollary}{6.10}\label{XI.6.10}
Let $X$ be a smooth $S$-group prescheme of finite presentation, with connected fibers, $Y$ a group prescheme of finite presentation over $S$, $i:Y\to X$ a monomorphism of $S$-group preschemes, thus making $Y$ a subgroup of $X$. Then $\prod_{X/S}Y/X$ is representable by a closed sub-prescheme of finite presentation of $S$. If $S$ is quasi-compact, denoting for every integer $n\ge0$ by $X_n$ the infinitesimal neighborhood of order $n$ of the identity section of $X$, and putting $Y_n=X_n\cap Y$, one has for large enough $n$: $\prod_{X/S}Y/X=\prod_{X_n/S}Y_n/X_n$.
\end{corollary}
The proof is essentially that of 6.9. Note that the identity sections of $X$ and $Y$ induce bijective immersions $S\to X_n$ and $S\to Y_n$, hence induce isomorphisms of $S_{\mathrm{red}}$ with $(X_n)_{\mathrm{red}}$ and $(Y_n)_{\mathrm{red}}$, which implies that the injection $Y_n\to X_n$ is proper, hence, being a monomorphism of finite presentation, is a closed immersion. Consequently, by VIII~6.4 already used, $\prod_{X_n/S}Y_n/X_n$ is representable by a closed sub-prescheme of $S$ of finite presentation over $S$, and it therefore remains to prove the last assertion of 6.10 in the case where one supposes in addition $S$ affine. One is again immediately reduced to the case where $S$ is noetherian, and one is reduced to proving that one then has $R=T$ (with the notation of the proof of 6.9), or again that $Y\supset X_n$ for every $n$ implies $Y=X$. Now the hypothesis implies that $i:Y\to X$ is étale at the points of the identity section of $Y$ over $S$, so $Y$ is smooth over $S$ at the points of the identity section, whence it follows that the open subset $Y'$ of points of $Y$ at which $Y$ is smooth over $S$ is an induced open subgroup of $Y$. Then $Y'\to X$ is an étale monomorphism by X~3.5, hence an open immersion; but since the fibers of $X$ are connected and every open subgroup of an algebraic group is also closed, it follows that this is a surjective open immersion, \ie an isomorphism. Thus $Y'=X$ and a fortiori $Y=X$, which completes the proof of 6.10.

Proceeding as in VIII~6.5, one concludes from 6.10:
\begin{corollary}{6.11}\label{XI.6.11}
Let $G$ be an $S$-group prescheme locally of finite type and quasi-separated over $S$, $H$ a group prescheme smooth of finite presentation over $S$ with connected fibers, $i:H\to G$ a monomorphism of $S$-groups, thus making $H$ a subgroup of $G$. Then:
\begin{enumerate}[label=\alph*)]
\item $\uCentr_G(H)$ and $\uNorm_G(H)$ are representable by closed sub-preschemes of $G$ of finite presentation over $G$, and likewise, for every monomorphism $j:K\to G$ of finite presentation of $S$-group preschemes, $\uTransp_G(H,K)$ is representable by a closed sub-prescheme of $G$ of finite presentation over $G$.
\item If $G$ is quasi-compact, in the various cases considered in a), there exists an integer $n\ge0$ such that (if $H_n$ denotes the infinitesimal neighborhood of order $n$ of the identity section of $H$) one has
\[
\begin{aligned}
\uCentr_G(H)&=\uCentr_G(H_n),\\
\uNorm_G(H)&=\uNorm_G(H_n),\\
\uTransp_G(H,K)&=\uTransp_G(H_n,K)=\uTransp_G(H_n,K_n).
\end{aligned}
\]
\end{enumerate}
\end{corollary}
One applies 6.10 to the group prescheme $X=H_G=H\times_S G$ over the base prescheme $G$, and to the group sub-prescheme $Y$ that is the inverse image of the diagonal subgroup of $(G\times_S G)_G$ over $G$ under a suitable homomorphism of $G$-groups from $X$ to $(G\times_S G)_G$ (in the case of $\Centr$), resp. the inverse image of $K_G$ under a suitable homomorphism of $G$-groups from $X$ to $G_G$ (in the case of $\Transp$). The case of $\Norm$ reduces to the transporter by taking $K=H$, the hypothesis on $G$ ensuring that $H\to G$ is of finite presentation (hence $Y\to X$ is of finite presentation); in the case of $\Centr$, the hypothesis made on $G$ ensures that the diagonal group of $G\times_S G$ is of finite presentation over $G\times_S G$, whence again the fact that $Y$ is of finite presentation over $X$.
\begin{remark}{6.12}\label{XI.6.12}
One can prove (using rather delicate results of EGA~VI{\renewcommand{\thefootnote}{*}\footnote{\Cf also J.~P.~Murre, \textit{Representation of unramified functors. Applications}, Sém.~Bourbaki \No294 (May 1965), th.~3 (p.~13).}\addtocounter{footnote}{-1}}) that if $X$ is a flat prescheme of finite presentation over $S$, which is proper over $S$ or has nonempty connected geometric fibers, then for every closed sub-prescheme $Y$ of $X$ of finite presentation over $S$, $\prod_{X/S}Y/X$ is representable by a closed sub-prescheme of $S$ of finite presentation over $S$. Likewise, if $X$ is a flat $S$-group prescheme of finite presentation with connected fibers, and $i:Y\to X$ a monomorphism of $S$-groups, with $Y$ an $S$-group prescheme of finite presentation, then $\prod_{X/S}Y/X$ is representable by a closed sub-prescheme of $S$ of finite presentation over $S$. In particular, 6.11 a) remains valid on replacing the hypothesis ``$H$ smooth over $S$'' by ``$H$ flat over $S$''.
\end{remark}
